\begin{longtable}{ |p{0.5cm}|p{3cm}|p{7cm}|p{12cm}|  }
	\caption{Tabel Referensi}\label{tab:referensi} \\
	\hline
	\multicolumn{1}{|c|}{\textbf{No}} & \multicolumn{1}{|c|}{\textbf{Sumber}} & \multicolumn{1}{|c|}{\textbf{Hasil Aplikasi}} & \multicolumn{1}{|c|}{\textbf{Perbedaan}} \\ \hline
	1 
 & Pengembangan Sistem Informasi Presensi Menggunakan Metode Waterfall.
 \newline Tri Wahyudi, Supriyanta Supriyanta, Husni Faqih.
 \newline Indonesian Journal On Software Engineering (IJSE)
 \newline ISSN : 2714-9935 
 & Sistem Informasi Presensi yang dikembangkan dengan metode pengembangan software Waterfall menghasilkan aplikasi yang berfungsi untuk menyediakan solusi terhadap masalah presensi karyawan dan siswa selama pandemi Covid-19. Aplikasi ini terdapat 2 pengguna yaitu karyawan dan admin. Karyawan dapat melakukan presensi serta melihat \textit{history} presensi, sedangkan admin dapat mengelola data karyawan, jam kerja karyawan, mengelola data presensi, dan mencetak rekap presensi. 
 & Pada jurnal sistem informasi presensi yang dikembangkan dengan metode pengembangan software waterfall, terdapat 2 \textit{role user} yaitu karyawan dan admin. Karyawan dapat melakukan presensi secara \textit{online} tapi data lokasi dan foto tidak dicatat, serta tidak terdapat sistem untuk melakukan perhitungan jam lembur.
 \newline 
  \newline
  Sedangkan pada aplikasi presensi \textit{online} berbasis \textit{web} pada PT Password Solusi Sistem menggunakan metode \textit{scrum}. Aplikasi ini terdapat 3 \textit{role user} yaitu \textit{manager} HRD, \textit{manager Infrastructure}, dan \textit{staff}. \textit{Staff} dapat melakukan presensi secara \textit{online} dan data lokasi serta foto akan dicatat oleh sistem. Selain itu, sistem akan melakukan perhitungan lembur jika \textit{staff} melakukan presensi diluar jam kerja atau pada hari libur.
 \\ \hline
 
 2 & APLIKASI PRESENSI ONLINE MENGGUNAKAN VALIDASI JARAK LOKASI PENGGUNA BERBASIS ANDROID (STUDY KASUS: TOKO YONIX)
 \newline Bayu Adytia Permana, Kisworo, Akhmad Jayadi.
 \newline Jurnal Informatika dan Rekayasa Perangkat Lunak (JATIKA)
 \newline ISSN : 2797-2011 
 & Aplikasi presensi \textit{online} menggunakan validasi jarak lokasi pengguna berbasis Android (study kasus: toko Yonix) berhasil dibangun berdasarkan model pada jurnal ini dan validasi jarak berhasil diterapkan menggunakan perhitungan antar lokasi yaitu pengguna dan toko, kemudian karyawan berhasil melakukan absensi secara tepat dan sesuai harapan. Terdapat 2 \textit{role} pada aplikasi ini. \textit{Role} karyawan dapat \textit{login}, melihat lokasi antara pengguna dan toko, absen kehadiran, dan informasi presensi karyawan. Sedangkan \textit{role} admin dapat login, melihat absen harian, dan melihat absen berdasarkan tanggal.
 & Pada jurnal aplikasi presensi \textit{online} menggunakan validasi jarak lokasi pengguna berbasis Android (study kasus: toko Yonix) dikembangkan berbasis android dan terdapat 2 \textit{role} yaitu karyawan dan admin. Karyawan hanya bisa \textit{login} dan melakukan presensi \textit{online} dan saat presensi juga menggunakan fitur GPS tapi hanya untuk melihat jarak antara pengguna dan toko saja dan tidak ada pengambilan foto. Selanjutnya pada \textit{role} admin dapat login, menambah data karyawan, dan mencari data absensi karyawan.
 \newline 
  \newline
  Sedangkan pada aplikasi presensi \textit{online} berbasis \textit{web} pada PT Password Solusi Sistem dikembangakan berbasis \textit{web} dan terdapat 3 \textit{role}. \textit{Role staff} dapat \textit{register/login}, melihat jadwal kerja, melakukan absensi dengan fitur GPS untuk mencatat lokasi pengguna dan mengambil foto pengguna serta sistem akan mencatat dan menghitung jam kerja lembur jika presensi dilakukan diluar jam kerja. Selanjutnya \textit{role manager} Infra yaitu dapat \textit{register/login}, dapat melakukan \textit{reject} pada data lembur yang tidak valid serta melihat data lembur \textit{staff}. Setelah itu, \textit{role manager} HRD dapat melakukan pengaktifan akun jika pengguna sudah \textit{register}, melihat data presensi dan lembur, serta mengelola data hari libur nasional dan jam kerja karyawan yang digunakan untuk mendeteksi jam lembur saat presensi.
 \\ \hline

  3 & Sistem Informasi Presensi Karyawan Berbasis Android (Studi Kasus: Asuransi Panin Dai-Ichi Life)
 \newline Tasya Rizki Ramadhini, Fenty Ariany, Akhmad Jayadi
 \newline Jurnal Teknologi dan Sistem Informasi (JTSI)
 \newline ISSN : 2746-3699
 & Sistem Informasi Presensi Karyawan Berbasis Android (Studi Kasus: Asuransi Panin Dai-Ichi Life) menggunakan implementasi GPS dengan metode geofence untuk mencari posisi devide saat karyawan melakukan presensi dan menggunakan metode Haversine Formula yang digunakan untuk mengetahui jarak radius lokasi agar dapat melakukan presensi. Pada aplikasi ini dirancang untuk pegawai saja yang dapat melakukan presensi dalam radius kantor dan pengambilan foto, terdapat menu form yang dapat digunakan oleh pegawai untuk izin atau sakit, terdapat menu untuk melihat history presensi, dan menu untuk melihat jadwal kerja.
 & Pada jurnal Sistem Informasi Presensi Karyawan Berbasis Android (Studi Kasus: Asuransi Panin Dai-Ichi Life) dikembangkan berbasis Android dan menggunakan metode \textit{waterfall}. Pada aplikasi ini menggunakan GPS hanya untuk membatasi radius saat melakukan presensi. Aplikasi ini hanya dirancang untuk \textit{role} pegawai saja yang dapat melakukan presensi, pengajuan izin atau sakit, melihat jadwal kerja, dan riwayat presensi.
 \newline 
  \newline
  Sedangkan pada aplikasi presensi \textit{online} berbasis \textit{web} pada PT Password Solusi Sistem dikembangkan berbasis \textit{web} dan menggunakan metode \textit{scrum}. Aplikasi ini menggunakan GPS untuk pencatatan lokasi saat presensi dan dikembangkan untuk 3 \textit{role} yaitu \textit{staff Infrastructure}, \textit{manager Infrastructure}, dan \textit{manager} HRD. Aplikasi ini juga dilengkapi sistem untuk melakukan pencatatan dan perhitungan lembur jika \textit{staff} melakukan presensi diluar jam kerja dan terdapat \textit{form} pengajuan lembur jika \textit{staff} lupa presensi saat lembur, akan tetapi aplikasi ini tidak ada \textit{form} untuk pengajuan izin atau sakit.
 \\ \hline

  4 & RANCANG BANGUN APLIKASI SISTEM INFORMASI ABSENSI KARYAWAN ONLINE
 \newline Rully Roosdianto, Ani Oktarini Sari, Arief Satriansyah
 \newline Jurnal Inti Nusa Mandiri
 \newline ISSN : 2685-807X
 & Rancang bangun aplikasi sistem informasi absensi karyawan online menghasilkan aplikasi yang dapat mempermudah karyawan dalam proses absensi, melakukan pencarian data absensi, menghitung seluruh rekap absensi karyawan, serta mengurangi resiko kehilangan dan kesalahan dalam pencatatan absensi pada CV Cahaya Toner. Aplikasi ini berbasis \textit{web} yang terdiri dari 2 \textit{role} yaitu \textit{admin} dan karyawan. \textit{Admin} dapat mengelola data karyawan, melihat rekap absensi, mengelola divisi, dan mengelola jam kerja karyawan. Sedangkan pada \textit{role} karyawan hanya dapat melakukan absensi dan melihat riwayat absensi.
 & Pada jurnal Rancang Bangun Aplikasi Sistem Informasi Absensi Karyawan Online dirancang berbasis \textit{web} menggunakan \textit{framework Codeigniter} dan menggunakan metode \textit{waterfall}. aplikasi ini hanya terdiri dari 2 \textit{role} yaitu \textit{admin} dan karyawan. Karyawan dapat melakukan absensi secara \textit{online} tapi saat proses tersebut tidak dilakukan pencatatan lokasi dan pengambilan foto.
 \newline 
  \newline
  Sedangkan pada aplikasi presensi \textit{online} berbasis \textit{web} pada PT Password Solusi Sistem dikembangkan berbasis \textit{web} dengan \textit{framwork Laravel} dan menggunakan metode \textit{scrum}. Aplikasi ini terdiri dari 3 \textit{role} yaitu \textit{staff Infrastructure}, \textit{manager Infrastructure}, dan \textit{manager} HRD. \textit{Staff Infrastructure} dapat melakukan presensi secara \textit{online} dengan tambahan teknik \textit{geolocation} yang dapat digunakan untuk mendapat titik koordinat \textit{staff} dan mencatat lokasi tersebut, aplikasi ini juga melakukan pengambilan foto saat presensi serta dilengkapi sistem untuk mendeteksi dan menghitung jam lembur jika \textit{staff} melakukan presensi diluar jam kerja. Data lembur yang masuk memiliki fitur \textit{reject} dari \textit{manager Infrastructure} atau penolakkan jika data lembur tidak \textit{valid}.
 \\ \hline

 
   5 & PENGEMBANGAN SISTEM ABSENSI ONLINE BERBASIS WEB MENGGUNAKAN MAPS JAVASRIPTS API
 \newline Gerlan Apriandy Manu, Yonly Adrianus Benufinit
 \newline Jurnal Pendidikan Teknologi Informasi (JUKANTI)
 \newline ISSN : 2621-1467
 & Pengembangan sistem absensi \textit{online} berbasis \textit{web} menggunakan \textit{maps javascripts API} menghasilkan aplikasi absensi \textit{online} yang dapat menyimpan \textit{latitude} dan \textit{longitude} pada \textit{user} yang melakukan absensi \textit{online} sehingga aplikasi ini dapat membantu manajemen absensi karyawan saat WFH (\textit{Work From Home}), yaitu melakukan absensi dimana pun dan kapanpun tanpa harus datang ke kantor menggunakan absensi \textit{fingerprint}. Aplikasi ini terdiri dari 2 \textit{role} yaitu pegawai dan petugas absensi. Pegawai dapat melakukan absensi secara \textit{online}, aplikasi ini menggunakan \textit{javascripts API} untuk mendapatkan lokasi dan mencatatnya, aplikasi juga mengambil foto selfie pegawai saat absensi, serta melihat riwayat absensi. Selanjutnya \textit{role} petugas absensi dapat melihat rekap absensi karyawan setiap bulan pada tanggal 20.
 & Pada jurnal Pengembangan sistem absensi \textit{online} berbasis \textit{web} menggunakan \textit{maps javascripts API} dikembangkan dengan \textit{PHPMaker} dan metode RAD (\textit{Rapid Application Development}). Aplikasi ini terdiri dari 2 \textit{role} yaitu pegawai dan petugas absensi. Aplikasi ini menggunakan API untuk mengambil dan mencatat lokasi karyawan dan foto pegawai saat presensi, namun tidak terdapat fitur jadwal kerja karyawan dan tidak bisa mendeteksi dan menghitung jumlah jam lembur pegawai jika presensi dilakukan pada hari libur atau diluar jam kerja.
 \newline 
  \newline
  Sedangkan pada aplikasi presensi \textit{online} berbasis \textit{web} pada PT Password Solusi Sistem dikembangkan berbasis \textit{web} dengan \textit{framwork Laravel} dan menggunakan metode \textit{scrum}. Aplikasi ini terdiri dari 3 \textit{role} yaitu \textit{staff Infrastructure}, \textit{manager Infrastructure}, dan \textit{manager} HRD. Aplikasi ini juga menggunakan bantuan API untuk mendapat lokasi \textit{staff} saat presensi namun dengan tambahan sistem untuk menghitung dan mencatat jam lembur \textit{staff} jika presensi dilakukan diluar \textit{shift} yang ditentukan. Data lembur yang masuk dapaat di-\textit{reject} oleh \textit{Manager Infrastructure} jika data lembur tersebut tidak \textit{valid}.
 \\ \hline

  6 & Perancangan Aplikasi Absensi Online Dengan Menggunakan Bahasa Pemrograman Kotlin
 \newline Arafat Febriandirza
 \newline Jurnal Pseudocode
 \newline ISSN : 2655-1845
 & Perancangan aplikasi absensi online dengan menggunakan bahasa pemrograman kotlin menghasilkan sistem yang membuat karyawan menjadi lebih disiplin, mengurangi kemungkinan adanya kecurangan saat melakukan absensi, meningkatkan efisiensi dan akurasi pencatatan data dan mengetahui lokasi karyawan saat absensi. Aplikasi ini memiliki 2 \textit{role} yaitu karyawan dan HRD karyawan. Karyawan dapat melakukan absensi \textit{online}, melihat data absensi, melakukan penugasan dan melihat data penugasan. Aplikasi ini dilengkapi GPS untuk mendapatkan alamat lokasi saat absensi.
 & Pada jurnal Perancangan aplikasi absensi online dengan menggunakan bahasa pemrograman kotlin dirancang berbasis Android dengan bahasa pemrograman \textit{kotlin} dan menggunakan metode \textit{agile}. Aplikasi ini dibagi menjadi 2 \textit{role} yaitu karyawan dan HRD Karyawan. Aplikasi ini sama-sama menggunakan GPS untuk mendapatkan lokasi \textit{user} namun aplikasi ini tidak terdapat pengambilan foto saat absensi dan tidak ada jadwal \textit{shift} karyawan dan sistem untuk mencatat dan menghitung jam lembur jika karyawan melakukan absensi diluar jam kerja atau pada hari libur.
 \newline 
  \newline
  Sedangkan pada aplikasi presensi \textit{online} berbasis \textit{web} pada PT Password Solusi Sistem dikembangkan berbasis \textit{web} dengan \textit{framwork Laravel} dan menggunakan metode \textit{scrum}. Aplikasi ini terdiri dari 3 \textit{role} yaitu \textit{staff Infrastructure}, \textit{manager Infrastructure}, dan \textit{manager} HRD. Aplikasi ini dilengkapi pengambilan foto \textit{user} saat melakukan presensi dan fitur tambahan untuk mendeteksi jam lembur berdasarkan \textit{shift} kerja \textit{staff}, jika \textit{staff} melakukan presensi diluar jam kerja maka akan masuk ke dalam data lembur. Data lembur dapat di-\textit{reject} oleh \textit{Manager Infrastructure} jika data lembur tersebut tidak \textit{valid}.
 \\ \hline

   7 & PERANCANGAN SISTEM INFORMASI ABSENSI ONLINE DETEKSI LOKASI BERBASIS WEB
 \newline Bryan Daniel Pesik, Penidas Fiodinggo Tanaem
 \newline Jurnal Mahasiswa Teknik Informatika (JATI)
 \newline ISSN : 2598-828X
 & Perancangan sistem informasi absensi \textit{online} deteksi lokasi berbasis \textit{web} menghasilkan aplikasi yang mempermudah kantor BPKAD untuk melakukan absensi dan merekap data laporan absensi karena data yang tersimpan didalam database. Aplikasi ini memiliki 2 \textit{role} yaitu pegawai dan \textit{admin}. Pegawai dapat melakukan absen pulang dan absen datang. Kemudian pada \textit{role admin} dapat merekap data absensi, mengatur jam kerja, melihat lokasi absen dan melihat data pegawai.
 & Pada jurnal Perancangan sistem informasi absensi \textit{online} deteksi lokasi berbasis \textit{web} dikembangkan dengan PHP dan \textit{framework Bootstrap}. Aplikasi ini memiliki 2 \textit{role} yaitu pegawai dan \textit{admin}. Pegawai dapat melakukan absensi tanpa fitur tambahan perhitungan lembur jika absensi dilakukan diluar jam kerja atau hari libur. Pada \textit{role admin} dapat mengatur jam kerja tapi untuk seluruh pegawai jadi kurang fleksibel.
 \newline 
 \newline
 Sedangkan pada aplikasi presensi \textit{online} berbasis \textit{web} pada PT Password Solusi Sistem dikembangkan berbasis \textit{web} dengan \textit{framwork Laravel} dan menggunakan metode \textit{scrum}. Aplikasi ini terdiri dari 3 \textit{role} yaitu \textit{staff Infrastructure}, \textit{manager Infrastructure}, dan \textit{manager} HRD. \textit{Staff} yang melakukan presensi diluar jam kerja atau hari libur akan masuk datanya ke data lembur. jika data lembur tidak \textit{valid} maka dapat di-\textit{reject} oleh \textit{Manager Infrastructure}. Kemudian pada \textit{role admin}, jadwal kerja dapat diatur untuk setiap \textit{staff} sehingga jam kerja setiap \textit{staff} dapat berbeda-beda atau fleksibel pengaturannya.
 \\ \hline

  8 & Perancangan Aplikasi Presensi Berbasis Web untuk
Karyawan Magang di PT. Sembilan Pilar Semesta
 \newline Marcel Alexandro Rumbang, Tony
 \newline Jurnal Ilmu Komputer dan Sistem Informasi (JIKSI)
 \newline ISSN : 2303-2529
 & Perancangan Aplikasi Presensi Berbasis Web untuk
Karyawan Magang di PT. Sembilan Pilar Semesta memudahkan proses presensi yang lebih praktis bagi karyawan magang dan memudahkan manajemen PT Sembilan Pilar Semesta dalam mengelola serta memantau kehadiran karyawan, berkat fitur pelacakan lokasi yang mengurangi risiko kecurangan. Aplikasi ini menggunakan \textit{framework Laravel}. Aplikasi ini memiliki 2 \textit{role} yaitu karyawan magang dan \textit{admin}. Karyawan magang dapat melakukan presensi datang maupun pulang dan melihat data presensi. Kemudian \textit{role admin} dapat melihat data presensi dan mengelola data karyawan.
 & Pada jurnal Perancangan Aplikasi Presensi Berbasis Web untuk
Karyawan Magang di PT. Sembilan Pilar Semesta dikembangkan dengan metode \textit{waterfall}. Aplikasi ini terdiri dari 2 \textit{role} karyawan magang dan admin. Aplikasi ini hanya mengambil lokasi saja tanpa melakukan pengambilan foto saat melakukan presensi dan tidak dapat melakukan presensi pulang ketika jam kerja belum sampai 8 jam. Aplikasi ini tidak dilengkapi fitur untuk menghitung lembur jika presensi dilakukan diluar jam kerja atau hari libur.
 \newline 
 \newline
 Sedangkan pada aplikasi presensi \textit{online} berbasis \textit{web} pada PT Password Solusi Sistem dikembangkan menggunakan metode \textit{scrum}. Aplikasi ini terdiri dari 3 \textit{role}. Aplikasi ini akan mengambil lokasi dan foto \textit{staff} saat presensi dan presensi pulang dapat dilakukan tanpa minimal total jam kerja. Aplikasi juga dilengkapi fitur untuk menghitung jam lembur jika presensi dilakukan diluar jam kerja. Data lembur yang masuk dapat di-\textit{reject} oleh \textit{Manager Infrastructure} jika data lembur tersebut tidak \textit{valid}. Pada \textit{role manager} HRD dapat mengatur jam kerja setiap \textit{staff} sehingga fleksibel yaitu dapat sama ataupun berbeda dengan \textit{staff} lainnya. \textit{manager} HRD juga dapat mengelola data hari libur kerja yang digunakan untuk mengecek lembur saat presensi bersamaan dengan jadwal kerja karyawan.
 \\ \hline
\end{longtable}


